\documentclass[12pt]{article}
\usepackage[utf8]{inputenc}
\usepackage[T1]{fontenc}
\usepackage{hyperref}
\title{Chinese, Huizhou Language Resources for K-12 Curriculum Development}
\author{Compiled by Claude Sonnet 5.5 (claude-sonnet-5-5) for Andr\'as Kornai}
\date{}
\begin{document}
\maketitle

\section{Introduction}
Huizhou Chinese (ISO 639-3 \texttt{czh}) is the Hui group of Chinese varieties in southern Anhui and adjacent areas, with several million speakers; nearly every county has its own variety. It should not be confused with the Huizhou dialect of Guangdong (Huicheng), a different variety of about 110,000 speakers discussed in some sources.

\section{Primary Educational Resources}
No textbooks, primers or school programmes for Huizhou Chinese were found in search results.

\section{Linguistic and Reference Materials}
\begin{itemize}
\item Wikipedia, Huizhou Chinese (sub-varieties Ji--She, Xiu--Yi, Qi--De): \url{https://en.wikipedia.org/wiki/Huizhou_Chinese}
\item Wikipedia, Huizhou dialect (the Guangdong variety, listed to avoid confusion): \url{https://en.wikipedia.org/wiki/Huizhou_dialect}
\item Forum thread on books about the Guangdong Huizhou dialect (mentions Liu Ruoyun's reference work; not about Anhui Hui): \url{https://chinalanguage.com/forums/viewtopic.php?p=21523}
\end{itemize}

\section{Religious and Community Resources}
None found.

\section{Digital and Audio Resources}
None found.

\section{Curriculum Development Context}
Little accessible material was found. Because the varieties differ greatly from county to county, any curriculum would first need to choose a variety and work with local speakers and Chinese-language dialectology literature (library catalogues and Chinese-language sources were not covered by these searches).
\end{document}